\section*{अध्याय \ref{ch:b3:many-electron-atoms} --- बहु-इलेक्ट्रॉन परमाणु और पद प्रतीक}

\begin{solution}{exo:b3:many-electron-atoms:1}
\[
\Psi = \frac{1}{\sqrt6}\begin{vmatrix}1s\alpha(1) & 1s\beta(1) & 2s\alpha(1)\\
1s\alpha(2) & 1s\beta(2) & 2s\alpha(2)\\ 1s\alpha(3) & 1s\beta(3) & 2s\alpha(3)
\end{vmatrix}.
\]
$1s$ केवल दो \omterm{def:b3:many-electron-atoms:spin-orbital}{चक्रण-कक्षक} देता है, $1s\alpha$ और $1s\beta$; तीसरा इलेक्ट्रॉन
उनमें से एक को दोहराता, जिससे दो स्तंभ बराबर हो जाते: सारणिक शून्य है।
\end{solution}

\begin{solution}{exo:b3:many-electron-atoms:2}
$\binom62 = 15$, $\binom63 = 20$, $\binom{10}{2} = 45$। $p^3$: ${}^4$S (4) +
${}^2$D (10) + ${}^2$P (6) = 20। $d^2$: ${}^3$F (21) + ${}^3$P (9) + ${}^1$G (9) +
${}^1$D (5) + ${}^1$S (1) = 45।
\end{solution}

\begin{solution}{exo:b3:many-electron-atoms:3}
N ($2p^3$, आधा भरा): \termsym{4}{S}{3/2}। O ($2p^4$, आधे से अधिक):
\termsym{3}{P}{2}। F ($2p^5$): \termsym{2}{P}{3/2}। \ce{Ti^2+} ($3d^2$):
\termsym{3}{F}{2}। \ce{Fe^3+} ($3d^5$, आधा भरा, सभी चक्रण समांतर, $L = 0$):
\termsym{6}{S}{5/2}।
\end{solution}

\begin{solution}{exo:b3:many-electron-atoms:4}
${}^2$D: $J = \frac52$ (6 अवस्थाएँ) और $\frac32$ (4); $6 + 4 = 10 = 5 \times 2$।
${}^3$P: $J = 2$ (5), 1 (3), 0 (1); $5 + 3 + 1 = 9 = 3 \times 3$।
\end{solution}

\begin{solution}{exo:b3:many-electron-atoms:5}
अंतराल $16.4$ ($J = 1 - 0$) और $27.0$ ($2 - 1$): अनुपात 2 के स्थान पर 1.65।
पहले से $A = 16.4/1 = \qty{16.4}{cm^{-1}}$, दूसरे से $27.0/2 = \qty{13.5}{cm^{-1}}$।
एक ही $A$ सटीक नहीं बैठता: स्तर अन्य पदों (${}^1$D, ${}^1$S) के स्तरों के साथ
थोड़े मिश्रित हैं, और \omterm{def:b3:many-electron-atoms:russell-saunders}{रसेल--सॉन्डर्स युग्मन} एक
सन्निकटन है।
\end{solution}

\begin{solution}{exo:b3:many-electron-atoms:6}
$\lambda = 10^7/\tilde\nu$: \qty{589.76}{nm} (D$_1$) और \qty{589.16}{nm}
(D$_2$)। अंतर $\qty{17.20}{cm^{-1}} = \qty{2.13}{meV}$। $E(\frac32) -
E(\frac12) = \frac32A$, अतः $A = \qty{11.47}{cm^{-1}}$।
\end{solution}

\begin{solution}{exo:b3:many-electron-atoms:7}
$3p \to 3s$: अनुमत ($\Delta l = -1$)। $3d \to 3s$: वर्जित ($\Delta l = -2$)।
$3d \to 3p$: अनुमत। $4s \to 3s$: वर्जित ($\Delta l = 0$, समता में कोई
परिवर्तन नहीं)। $4s \to 3p$: अनुमत।
\end{solution}

\begin{solution}{exo:b3:many-electron-atoms:8}
${}^3$P का माध्य: $(5 \times 169086.76 + 3 \times 169086.84 + 169087.83)/9 =
\qty{169086.91}{cm^{-1}}$। ${}^1$P तक का अंतराल: \qty{2047.98}{cm^{-1}} $=
\qty{0.254}{eV}$, अतः $K(1s,2p) = \qty{0.127}{eV}$, जो
$K(1s,2s)$ से तीन गुना छोटा है। \omterm{def:b3:many-electron-atoms:exchange}{विनिमय समाकल} दोनों कक्षकों के घनत्वों के
अतिव्यापन को मापता है; $2s$ कक्षक $1s$ क्षेत्र में भेदन करता है (नाभिक के पास
उसका घनत्व होता है), $2p$ कक्षक नाभिक पर लुप्त होता है और कम अतिव्यापन करता है।
\end{solution}

\begin{solution}{exo:b3:many-electron-atoms:9}
सबसे बड़े $M_L = 4$ के लिए दोनों इलेक्ट्रॉन $m_l = 2$ में चाहिए, विपरीत चक्रणों के साथ: एक
एकल, ${}^1$G (9 अवस्थाएँ)। अगला, $M_L = 3$, $m_l = 2, 1$ से आता है, जहाँ
$M_S = 1$ संभव है: ${}^3$F (21)। शेष $M_L = 2$ केवल $M_S = 0$ के साथ:
${}^1$D (5)। शेष $M_L = 1$, $M_S = 1$ के साथ: ${}^3$P (9)। एक अवस्था बचती है,
$M_L = M_S = 0$: ${}^1$S। कुल 45।
\end{solution}

\begin{solution}{exo:b3:many-electron-atoms:10}
$\mathbf r_1 = \mathbf r_2 = \mathbf r$ पर त्रिक फलन
$\frac{1}{\sqrt2}[1s(\mathbf r)2s(\mathbf r) - 2s(\mathbf r)1s(\mathbf r)] = 0$ है;
एकल फलन $\sqrt2\,1s(\mathbf r)2s(\mathbf r) \ne 0$ है। त्रिक में
इलेक्ट्रॉन कभी एक ही बिंदु पर नहीं पाए जाते और औसतन अधिक दूर
होते हैं, अतः उनका प्रतिकर्षण कम है: $E_+ - E_- = 2K > 0$।
\end{solution}

\begin{solution}{exo:b3:many-electron-atoms:11}
$d^8$ आधे से अधिक भरा है: युग्मन प्रतिलोमित है, $J = L + S = 4$
सबसे नीचे। अंतराल: $E(3) - E(4) = 1360.7$ और $E(2) - E(3) = 908.9$; अनुपात
1.50, लांडे के $4/3 = 1.33$ के विरुद्ध। $|A| = 1360.7/4 = \qty{340}{cm^{-1}}$
($A < 0$)। इस भारी आयन में \omterm{def:b3:many-electron-atoms:spin-orbit}{चक्रण--कक्षक युग्मन} कार्बन
(\qty{16}{cm^{-1}}) की तुलना में कहीं अधिक है।
\end{solution}

\begin{solution}{exo:b3:many-electron-atoms:12}
$\sum_J(2J+1)[J(J+1) - L(L+1) - S(S+1)] = 0$ ($(2L+1)(2S+1)$ अवस्थाएँ
किसी भी आधार में गिनी जा सकती हैं, और अयुग्मित आधार पर $\sum M_LM_S = 0$)।
${}^3$P के लिए: $1(0 - 4) + 3(2 - 4) + 5(6 - 4) = -4 - 6 + 10 = 0$। कार्बन का
भारित माध्य $(0 + 3 \times 16.42 + 5 \times 43.41)/9 = \qty{29.6}{cm^{-1}}$ है:
वह ऊर्जा जो पद की \omterm{def:b3:many-electron-atoms:spin-orbit}{चक्रण--कक्षक युग्मन} के बिना होती।
\end{solution}

\begin{solution}{pb:b3:many-electron-atoms:1}
\textbf{1.} $1s\alpha$, $1s\beta$, $2s\alpha$, $2s\beta$।
\textbf{2.} $|1s\alpha\,2s\alpha|$, $|1s\alpha\,2s\beta|$, $|1s\beta\,2s\alpha|$,
$|1s\beta\,2s\beta|$।
\textbf{3.} $|1s\alpha\,2s\alpha| = \frac{1}{\sqrt2}[1s(1)2s(2) - 2s(1)1s(2)]\,
\alpha(1)\alpha(2)$, $2\times2$ सारणिक का प्रसार करके।
\textbf{4.} उनका योग $\frac{1}{\sqrt2}[1s(1)2s(2) - 2s(1)1s(2)]\cdot
\frac{1}{\sqrt2}[\alpha(1)\beta(2) + \beta(1)\alpha(2)]$ है ($\sqrt2$ गुना, फिर
प्रसामान्यीकृत); उनका अंतर $\frac{1}{\sqrt2}[1s(1)2s(2) +
2s(1)1s(2)]\cdot\frac{1}{\sqrt2}[\alpha(1)\beta(2) - \beta(1)\alpha(2)]$ देता है।
\textbf{5.} सममित आकाशी फलन प्रतिसममित चक्रण फलन के साथ ($S = 0$);
प्रतिसममित आकाशी फलन तीनों सममित चक्रण फलनों के साथ ($S = 1$)।
\textbf{6.} $2S + 1 = 1$ और 3: समान ऊर्जा की एक और तीन अवस्थाएँ।
\textbf{7.} $[\ce{Ar}]\,3d^2$; $\binom{10}{2} = 45$ सारणिक।
\textbf{8.} अधिकतम $S = 1$, अधिकतम $L = 2 + 1 = 3$, आधे से कम भरा:
\termsym{3}{F}{2}, आँकड़ों में 0 पर स्थित स्तर।
\textbf{9.} $21 + 9 + 9 + 5 + 1 = 45$।
\textbf{10.} सामान्य ($J = 2$ सबसे नीचे)। अंतराल 184.9 और 235.5: अनुपात 1.27,
$4/3 = 1.33$ के विरुद्ध; नियम 5\,\% के भीतर लागू होता है।
\textbf{11.} ${}^3$F: $(5 \times 0 + 7 \times 184.9 + 9 \times 420.4)/21 =
\qty{241.8}{cm^{-1}}$; ${}^3$P: $(10538.4 + 3 \times 10603.6 + 5 \times
10721.2)/9 = \qty{10661.7}{cm^{-1}}$।
\textbf{12.} $15B = \qty{10419.9}{cm^{-1}}$, $B = \qty{695}{cm^{-1}}$।
\textbf{13.} $\Delta S = 0$।
\textbf{14.} $\Delta S = 1$ वर्जित है, और $1s2s \to 1s^2$ में $\Delta l = 0$
(समता में कोई परिवर्तन नहीं; ${}^3$S$_1 \to {}^1$S$_0$ में दो S
अवस्थाओं के बीच $\Delta L = 0$ भी है)। यह अवस्था मितस्थायी है: वह सामान्य उत्तेजित
अवस्थाओं से कहीं अधिक देर तक जीवित रहती है।
\textbf{15.} $169086.91 - 159855.97 = \qty{9230.94}{cm^{-1}}$: \qty{1083.3}{nm},
निकट अवरक्त में।
\textbf{16.} $171134.89 - 166277.44 = \qty{4857.45}{cm^{-1}}$: \qty{2058.7}{nm}।
\textbf{17.} $\qty{171134.89}{cm^{-1}}$: \qty{58.43}{nm}, निर्वात
पराबैंगनी में (वायु द्वारा अवशोषित)।
\textbf{18.} हर कुल की अपनी रेखाएँ और अपनी निम्नतम अवस्था है, और
दोनों कभी प्रकाश द्वारा नहीं जुड़ते: वे दो स्पेक्ट्रमों वाली दो गैसों की तरह व्यवहार करते हैं।
\textbf{19.} $E_\pm = E_{1s} + E_{2s} + J \pm K$ (एकल $+$, त्रिक $-$)।
\textbf{20.} त्रिक: इलेक्ट्रॉनों के मिलने पर उसका आकाशी फलन लुप्त हो जाता है,
अतः वे कम प्रतिकर्षण करते हैं।
\textbf{21.} \qty{6421.47}{cm^{-1}}, $6421.47/8065.54 = \qty{0.796}{eV}$।
\textbf{22.} अंतराल $\qty{2047.98}{cm^{-1}} = \qty{0.254}{eV}$, अतः $K(1s,2p) =
\qty{0.127}{eV}$।
\textbf{23.} $2s$, $1s$ क्षेत्र में भेदन करता है और $2p$ की तुलना में उससे
अधिक अतिव्यापन करता है।
\textbf{24.} हाँ: दोनों विन्यासों में त्रिक (बड़ा $S$) सबसे नीचे है।
\textbf{25.} $K = \frac12 \times \qty{0.796}{eV}$: \textbf{हीलियम का \omterm{def:b3:many-electron-atoms:exchange}{विनिमय
समाकल} $K(1s,2s)$ \qty{0.398}{eV} है}, जो एकल--त्रिक अंतराल के आधे के रूप में
मापा गया।
\end{solution}
